\chapter{比例与数据}\label{ch:g6:propdata}

若 3 本笔记本售价 $6$ 欧元，6 本笔记本售价 $12$ 欧元：本数翻倍，
价格也翻倍。
这样变化的量称为\emph{\omterm{def:g6:propdata:def}{成正比例}}的量——这是数学中最有用的思想之一，
并在\cref{ch:g7:prop}中进一步展开。本章最后
阅读和绘制简单的数据统计图。

\section{成正比例的量}

\begin{definition}[成正比例]\label{def:g6:propdata:def}
若两个量\emph{成正比例}\index{成正比例}，则其中一个量的值
总是由另一个量的值乘
\emph{同一个}数得到，这个数称为\emph{比例
系数}。
\end{definition}

\begin{example}\label{ex:g6:propdata:table}
笔记本每本 $2$ 欧元：
\begin{center}
\renewcommand{\arraystretch}{1.3}
\begin{tabular}{l|cccc}
笔记本数 & $3$ & $5$ & $8$ & $12$ \\
\hline
价格（欧元） & $6$ & $10$ & $16$ & $24$ \\
\end{tabular}
\end{center}
每个价格都是笔记本数 $\times 2$，系数为
$2$（即单价）。快速检验一张表是否\omterm{def:g6:propdata:def}{成正比例}：
所有“列商”$\frac{6}{3}, \frac{10}{5}, \frac{16}{8},
\frac{24}{12}$ 必须相等——这里它们都等于 $2$。
\end{example}

\begin{example}[反例]\label{ex:g6:propdata:nonexample}
年龄与身高不\omterm{def:g6:propdata:def}{成正比例}：$12$ 岁的孩子
身高并不是 $6$ 岁孩子的两倍。用数值检验：$1.50$ m 对应 $12$ 岁，
$1.15$ m 对应 $6$ 岁，商分别为 $\frac{1.50}{12} = 0.125$ 和
$\frac{1.15}{6} \approx 0.19$，两者不相等。
\end{example}

\begin{method}[补全比例表]\label{met:g6:propdata:complete}
可任选以下三种方法：
\begin{enumerate}
  \item \emph{系数法}：由一个完整的列求出乘数，
        再将其用于其他各列；
  \item \emph{列法}：一列可以乘以某个数，或把两列相加，
        得到新的一列；
  \item \emph{归一法}：先求出\emph{一个}物品的值，
        再乘以所求的数量。
\end{enumerate}
\end{method}

\begin{example}\label{ex:g6:propdata:methods}
5 个相同的杯子共 $15$ 欧元，8 个这样的杯子共多少欧元？

\emph{归一法}：一个杯子的价格为 $15 \div 5 = 3$ 欧元，所以 8 个杯子
共 $8 \times 3 = 24$ 欧元。

\emph{列法}：$8 = 5 + 3$；3 个杯子的价格是 $\frac{3}{5}$ 的 $15$，
即 $9$ 欧元，因此 8 个杯子共 $15 + 9 = 24$ 欧元。两种方法结果一致，
这是必然的。
\end{example}

\section{百分数}

\begin{definition}[百分数]\label{def:g6:propdata:percent}
“一个量的 $t\,\%$”表示该量的 $\frac{t}{100}$：
例如 $25\,\%$ 的 $60$ 为 $\frac{25}{100} \times 60 = 15$。百\omterm{def:g6:fractions:def}{分数}是
系数为 $\frac{t}{100}$ 的\omterm{def:g6:propdata:def}{成正比例}关系。
\end{definition}

\begin{example}\label{ex:g6:propdata:percent}
某班有 $30$ 名学生，其中女生占 $40\,\%$。逐步求女生人数：
$40\,\%$ 表示 $\frac{40}{100} = 0.4$，于是
则女生有 $0.4 \times 30 = 12$ 人。一些便于记忆的标准值：
$50\,\%$ 是\omterm{def:g3:division:half}{一半}，$25\,\%$ 是\omterm{def:g3:division:half}{四分之一}，$10\,\%$ 是十分之一——
因此 $10\,\%$ 的 $30$ 是 $3$，而 $40\,\% = 4 \times 10\,\%$，
$4 \times 3 = 12$：结果相同，可以心算得出。
\end{example}

\section{读图与绘制统计图}

\begin{method}[读表或读图]\label{met:g6:propdata:read}
无论哪种图示——表格、条形图或折线图：
\begin{enumerate}
  \item 先读标题，以及两个坐标轴的单位；
  \item 回答问题时，先找到对应的条形（柱），
        再对照刻度仔细读数；
  \item 比较时按条形的高矮比较——但要检查坐标轴是否从 $0$ 开始，
        否则图形可能产生误导。
\end{enumerate}
\end{method}

\begin{example}\label{ex:g6:propdata:chart}
下面的条形图给出了一周内学校图书馆
借出的图书数量。

\begin{omfigure}
\begin{tikzpicture}
\begin{axis}[
  width=8.6cm, height=5.2cm,
  ybar, bar width=16pt,
  symbolic x coords={Mon,Tue,Wed,Thu,Fri},
  xtick=data, xticklabels={周一,周二,周三,周四,周五},
  ymin=0, ymax=22,
  ytick={5,10,15,20},
  ylabel={借出图书数},
  tick label style={font=\small},
  label style={font=\small}]
  \addplot[fill=omDef!30, draw=omDef] coordinates
    {(Mon,12) (Tue,8) (Wed,17) (Thu,6) (Fri,15)};
\end{axis}
\end{tikzpicture}

{\small 条形图：每天一个条形，高度等于频数。}
\end{omfigure}

读图：借书最多的是周三（$17$ 本）。周二和
周四合计借出 $8 + 6 = 14$ 本，少于周五
一天的 $15$ 本。一周总数为：
$12 + 8 + 17 + 6 + 15 = 58$ 本。
\end{example}

\begin{omfigure}
\begin{tikzpicture}
\begin{axis}[omaxis,
  width=8.2cm, height=5.4cm,
  xmin=0, xmax=9.4, ymin=0, ymax=28,
  xlabel={笔记本数}, ylabel={价格（欧元）},
  xtick={1,...,9}, ytick={4,8,12,16,20,24}]
  \addplot[omDef, thick, domain=0:9] {3*x};
  \addplot[omDef, only marks, mark=*, mark size=1.5pt] coordinates
    {(1,3) (2,6) (3,9) (4,12) (5,15) (6,18) (7,21) (8,24)};
  \draw[dashed, omThm] (axis cs:6,0) -- (axis cs:6,18)
    -- (axis cs:0,18);
\end{axis}
\end{tikzpicture}

{\small \omterm{def:g6:propdata:def}{成正比例}的量，其图像很容易辨认：各点
落在一条过原点的直线上（此处为笔记本，每本
$3$ 欧元；虚线读出 $6$ 本笔记本的价格为 $18$
欧元）。}
\end{omfigure}

\section{练习}

\begin{exercise}[$\star$]\label{exo:g6:propdata:1}
判断下表是否\omterm{def:g6:propdata:def}{成正比例}？请通过计算商来说明理由。
\begin{center}
\renewcommand{\arraystretch}{1.2}
\begin{tabular}{l|ccc}
苹果（千克） & $2$ & $3$ & $5$ \\
\hline
价格（欧元） & $5$ & $7.5$ & $12.5$ \\
\end{tabular}
\qquad
\begin{tabular}{l|ccc}
年龄（岁） & $2$ & $4$ & $8$ \\
\hline
身高（厘米） & $85$ & $103$ & $130$ \\
\end{tabular}
\end{center}
\end{exercise}

\begin{exercise}[$\star$]\label{exo:g6:propdata:2}
四个羊角面包售价 $4.80$ 欧元。求一个羊角面包的价格，
再求七个羊角面包的价格。
\end{exercise}

\begin{exercise}[$\star$]\label{exo:g6:propdata:3}
补全比例表（先求比例系数！）：
\begin{center}
\renewcommand{\arraystretch}{1.2}
\begin{tabular}{l|cccc}
燃油（升） & $10$ & $25$ & $?$ & $60$ \\
\hline
价格（欧元） & $18$ & $?$ & $72$ & $?$ \\
\end{tabular}
\end{center}
\end{exercise}

\begin{exercise}[$\star$]\label{exo:g6:propdata:4}
一辆汽车耗油 $6$ L 可行驶 $100$ km。行驶 $250$ km 需要多少燃油？
行驶 $350$ km 呢？用 $27$ L 燃油能行驶多远？
\end{exercise}

\begin{exercise}[$\star$]\label{exo:g6:propdata:5}
利用\cref{ex:g6:propdata:percent}中的基准数，口算：
$50\,\%$ 的 $84$；\quad $25\,\%$ 的 $200$；\quad $10\,\%$ 的 $63$；
\quad $30\,\%$ 的 $70$。
\end{exercise}

\begin{exercise}[$\star$]\label{exo:g6:propdata:6}
某学校有 $450$ 名学生，其中 $60\,\%$ 在食堂用餐。在食堂用餐的
有多少名学生？不在食堂用餐的有多少名？
\end{exercise}

\begin{exercise}[$\star$]\label{exo:g6:propdata:7}
利用\cref{ex:g6:propdata:chart}中的条形图，哪几天
借出的图书少于 $10$ 本？周三借出的图书
比周四多多少本？
\end{exercise}

\begin{exercise}[$\star$]\label{exo:g6:propdata:8}
周一至周五中午的气温分别为 $14$、$16$、$13$、
$17$、$20$ 度。请画出这些数据的条形图（选择合适的
刻度），并读出气温最高的一天。
\end{exercise}

\begin{exercise}[$\star\star$]\label{exo:g6:propdata:9}
某菜谱供 $6$ 人食用，需要 $450$ g 面粉和 $3$ 个鸡蛋。请将其调整
为供 $10$ 人食用的用量。（关于鸡蛋，先想一想，
鸡蛋的个数能是\omterm{def:g6:decimals:places}{小数}吗？）
\end{exercise}

\begin{exercise}[$\star\star$]\label{exo:g6:propdata:10}
一位女骑车人以匀速骑行，行驶 $24$ km 用时 $1$ 小时。
\begin{enumerate}
  \item 她骑行 $2$ 小时能走多远？骑行半小时呢？骑行
        $1$ 小时 $30$ 分钟呢？
  \item 她骑行 $60$ km 需要多长时间？
\end{enumerate}
\end{exercise}

\begin{exercise}[$\star\star$]\label{exo:g6:propdata:11}
商店提供“全场 $20\,\%$ 折扣”。求下列商品的折扣金额和
折后价：一个球 $15$ 欧元；一个球拍 $40$ 欧元。
折后价是否与原价\omterm{def:g6:propdata:def}{成正比例}？比例系数是多少？
\end{exercise}

\begin{exercise}[$\star\star\star$]\label{exo:g6:propdata:12}
两根高度相同的蜡烛同时点燃。粗的一根
以各自恒定的速率燃烧，$6$ 小时恰好烧完，细的一根需 $4$ 小时。
经过多长时间后，细蜡烛恰好是粗蜡烛高度的\omterm{def:g3:division:half}{一半}？
（用\omterm{def:g6:fractions:def}{分数}表示 $t$ 小时后的
剩余高度占初始高度的比例。）
\end{exercise}

\section{问题：地图、平面图与如何用图表“说谎”}

\begin{problem}[{周末问题——地图的比例尺是一种比例关系：长度按系数变化，
面积按系数的平方变化，
而画错的图表会误导人的视觉}]
\label{pb:g6:propdata:1}
每张地图、平面图和模型都遵循同一条规律：实际长度
与图上长度\emph{\omterm{def:g6:propdata:def}{成正比例}}
（见\cref{def:g6:propdata:def}）。这个比例系数有一个著名的名字，
叫作\emph{比例尺}\index{比例尺}，它隐藏着两个陷阱，
本题有意设置这两个陷阱：\omterm{def:g6:measure:area}{面积}\emph{并不}按比例尺变化，
而百\omterm{def:g6:fractions:def}{分数}和图表一样，完全取决于它们
所比较的基准。

\medskip
\noindent\textbf{第一部分——读地图。}
一张徒步地图的比例尺为 $1:100\,000$：地图上每
厘米代表地面上 $100\,000$ 厘米。
\begin{enumerate}
  \item 将 $100\,000$ cm 换算为千米。这张地图上 $1$ cm
        代表多少 km？
  \item 地图上两个村庄相距 $7.5$ cm；湖岸是一条
        $12$ cm 长的曲线。求它们的实际距离。
  \item 一条笔直的罗马古道长 $20$ km。它在地图上
        有多长？
  \item 第二张地图标明“$1$ cm 代表 $5$ km”。把它的比例尺写成
        $1: \,?$ 的形式。对同一地区，两张地图中
        哪一张显示的细节更多？
  \item 为徒步地图列一张小表（地图上距离分别为 $1$、$2$、$7.5$ cm，
        对应实际距离），并说明为什么比例尺恰好是
        一种\omterm{def:g6:propdata:def}{成正比例}关系，依据
        \cref{def:g6:propdata:def}的含义。把图上厘米数换算为千米数的
        系数是多少？
\end{enumerate}

\medskip
\noindent\textbf{第二部分——卓伊的卧室平面图。}
卓伊画了她卧室的平面图——一个长 $4$ m、宽
$3$ m 的长方形——比例尺为 $1:50$。
\begin{enumerate}[resume]
  \item 房间在平面图上的尺寸是多少？
  \item 她的床长 $2$ m、宽 $0.9$ m，门洞
        宽 $80$ cm。求这三个尺寸在平面图上的长度。
  \item 现在来看陷阱。先求房间的实际\omterm{def:g6:measure:area}{面积}（单位为
        m$^2$），再求平面图的\omterm{def:g6:measure:area}{面积}（单位为 cm$^2$）。
        把实际\omterm{def:g6:measure:area}{面积}换算为 cm$^2$
        （见\cref{def:g6:measure:area}），再除以平面图\omterm{def:g6:measure:area}{面积}：
        商是否为 $50$？
  \item 解释你得到的数：在 $1:50$ 的平面图上，纸上每平方厘米
        代表实际中一个长为 $50$ cm、宽为 $50$ cm 的正方形。
        这相当于多少实际的 cm$^2$？请总结这一规律：
        当长度除以 $50$ 时，\omterm{def:g6:measure:area}{面积}就要除以
        \,\dots
  \item 圆形地毯在平面图上覆盖 $6$ cm$^2$。它
        实际覆盖的\omterm{def:g6:measure:area}{面积}是多少？先以 cm$^2$，再以 m$^2$ 表示。
\end{enumerate}

\medskip
\noindent\textbf{第三部分——如何用图表
（和百\omterm{def:g6:fractions:def}{分数}）“说谎”。}
\begin{enumerate}[resume]
  \item 一家商店周六的销售额为 $105$ 欧元，周日为 $110$
        欧元。经理画了一张条形图，其纵轴
        从 $100$ 开始：周六的条形上升 $5$ 个小格，
        周日的条形上升 $10$ 个小格。这张图给人什么暗示，
        实际情况又如何？（以百\omterm{def:g6:fractions:def}{分数}表示周日相对周六的
        增长量，精确到整数百分比。）
        经理忽略了下列哪一条警示（见
        \cref{met:g6:propdata:read}）？
  \item 重新画一张（或描述一张）诚实的图表，纵轴
        从 $0$ 开始：此时两个条形如何比较？
  \item 一个收藏册中的邮票从 $40$ 枚增加到 $60$ 枚，求
        增长量占初始值的百\omterm{def:g6:fractions:def}{分数}。之后它从 $60$ 枚减少回
        $40$ 枚，求减少量占\emph{其}初始值的百\omterm{def:g6:fractions:def}{分数}。
        同样是 $20$ 枚邮票，
        为什么两个答案不同？
  \item 换季促销：一件 $100$ 欧元的外套先打“$20\,\%$
        折扣”，结账时又加一个“折后价再打 $10\,\%$ 折扣”。
        请逐步计算最终价格。
        总折扣是否为 $30\,\%$？请向顾客解释它
        实际是多少。
  \item 结尾，回到第 4 题的地图（$1$ cm 代表 $5$ km）：一片森林
        在该地图上占 $3$ cm$^2$。它的实际\omterm{def:g6:measure:area}{面积}是多少
        km$^2$？最后以本题的整体规律作结：
        长度遵循比例尺，\omterm{def:g6:measure:area}{面积}遵循它的
        \,\dots
\end{enumerate}
\end{problem}
